\subsection{A passing requirement still needs the intended native request}
\label{sec:native-requirement-link}

Consider a regulatory assessment for the synthetic bank's solo legal book.
The supplied assessment date is 31 March 2026, the period starts on 1 January,
and its exclusive endpoint is 1 April. The accounting state is also dated
31 March. Selected requirements can verify that the period is nonempty, that
the assessment lies within it, and that the state and assessment dates agree.
These statements remain true if the request names an obsolete accounting head
or gives the operation a zero execution budget. The native request must still
be refused in those cases.

An input binding therefore needs an explicit connection to the operation being
requested. The proposed \code{NativeRequirementLink/1} supplies a finite example
for \code{bank.regulatory.assess@1}. It names the exact request, submission,
profile, operation signature, assessment basis and context. It also binds the
requirement manifest and input-plan digests. The selected assessment and period
dates come from the basis record; the state date comes from the accounting
state referenced by that basis. Each selected scalar keeps its pointer, state
role and full context, including bank, book, perimeter, currency, scale,
accounting basis and jurisdiction. A matching currency alone cannot justify
substituting another book or regulatory perimeter.

Two different content identities are involved. The native record digest includes
its schema identifier, whereas the binder's raw digest identifies the supplied
serialized bytes. The submission contains the complete canonical request
record. The basis and state inputs are explicitly serialized from their retained
native records; they are not original filing bytes. The link checks both kinds
of identity and their correspondence. Replacing a signature, selecting a
different record, weakening the selected definition or retaining a stale digest
causes refusal rather than an implicit conversion.

The existing binder and original version-selected evaluator still perform the
input and predicate work. The link adds no predicate language or dispatcher.
Its success means that these particular selected requirements are tied to this
particular native request. The native method must separately establish its
last-day condition, permitted scenario window, evidence availability and
ingestion times, source applicability, regulatory eligibility, current view and
accounting head, environment and resource budget. Request identity and replay
also remain part of the native operation. None of these duties disappears when
the selected date requirements succeed.

The retained comparison makes the distinction observable. Across 271 regulatory
exchanges, the native results include ten completed assessments, 75 replays,
50 unknown acknowledgements, 96 refusals and 40 partial assessments. Selected
date requirements pass in 95 of the 96 refused exchanges and in every partial
and unknown-acknowledgement case. An unknown acknowledgement leaves publication
unresolved; it cannot be interpreted as a failed assessment with no effects.
The original native response checks preserve those meanings. The three
requirement versions produce 813 matching text and typed assessments, with
independent date comparisons and malformed-input checks.

Missing evidence retains its own consequence. If a declared date is removed
from an explicitly derived input record, the binder reports that date unknown.
The derived record cannot then be represented as the unchanged native record.
An unavailable selected capability remains unsupported, and a known date
contradiction survives an unrelated unknown. Execution, source admission,
accounting admission and full method eligibility remain false in these link
results.

This example supplies a checked connection for one retained operation. The
portable corporate purchase definition, whose preconditions are named strings,
does not thereby become a definition of the bank assessment. General method
authoring still requires explicit correspondences for each supported method,
including its state and effect semantics. The checked source and result records
are retained under
\path{data/accounting-master/common-method-input-boundary-20261026/}.
